\documentclass[11pt,a4paper]{article}
\usepackage[margin=27mm]{geometry}
\usepackage{amsmath,amsthm,amssymb,fontspec,unicode-math,hyperref}
\usepackage{longtable,booktabs,array,calc,fvextra}
\setmainfont{STIXGeneral}
\setmathfont{STIX Two Math}
\setmonofont{Menlo}
\hypersetup{hidelinks}
\newtheorem{theorem}{Theorem}
\newtheorem{proposition}[theorem]{Proposition}
\newtheorem{lemma}[theorem]{Lemma}
\newtheorem{corollary}[theorem]{Corollary}
\theoremstyle{remark}\newtheorem{remark}[theorem]{Remark}
\newcommand{\gen}[1]{\langle #1\rangle}
\newcommand{\F}{\mathbb F}
\newcommand{\Z}{\mathbb Z}
\newcommand{\Aut}{\operatorname{Aut}}
\newcommand{\Sym}{\operatorname{Sym}}
\setlength{\emergencystretch}{2em}
\providecommand{\tightlist}{\setlength{\itemsep}{0pt}\setlength{\parskip}{0pt}}

\title{Halving generalized cubes}
\author{Human Compression Project}
\date{}
\begin{document}
\maketitle
\begin{abstract}
We apply the common graph-halving theorem to Danzer's generalized
cubes. The automorphism group is the augmentation module of the vertex
set extended by the original polytope group. This description recovers
the polygon parameters and all locally spherical cases in the stated
construction. Explicit Gray coordinates identify both the original and
halved toroidal face posets, distinguished generators and translation
lattices. These recover the principal conclusions of Montero and
Ivi\'c Weiss for their first arXiv edition, with its auxiliary rank
qualification and exclusions stated separately.
\end{abstract}
\section{The common halving theorem and source conventions}
The common paper [U] proves the endpoint intersection criterion
C.Endpoint and realizes the operation geometrically in C.Graph and
C.Halving. In particular, for a regular lattice polytope, replacing an
endpoint marker $a$ by $aba$ gives a regular hypertope. Its group is
normal of index at most two; the index is two exactly when the relevant
vertex graph is bipartite. If the old endpoint order is $m<\infty$, the
new order is $m/\gcd(m,2)$, and infinite order stays infinite.
The other labels incident with $b$ are duplicated. Reversing the types
gives the source's convention at the final endpoint.

These are [MW20, Theorem 3.1, Proposition 3.2 and Corollary 3.3], with
the same lattice hypothesis. The term \emph{non-degenerate} is not
replaced by an assumption about a diagram alone. The source's
Lemma 2.4 follows directly because a subgroup generated by disjoint
marked factors is their product and intersections are coordinatewise.
Its quoted Proposition 2.1 is background; the incidence recognition
needed from its quoted Theorem 2.2 is proved as UC.Geometry in [U].
The quoted Proposition 2.3 is not needed.

\section{Parity modules and locally spherical
geometries}\label{parity-modules-and-locally-spherical-geometries}

Retain the following Danzer construction. If K is a finite regular
lattice polytope with m vertices, \(2^K\) has binary vertex set
\(V=F_{2}^m\). A face associated with a face F of K and \(x\in V\)
allows exactly the coordinates indexed by the vertices of F to vary.
These coordinate subcubes, ordered by inclusion with their improper
faces, form a regular lattice polytope of one higher rank, with vertex
figures K and automorphism group

\[
\operatorname{Aut}(2^K)=V\rtimes\operatorname{Aut}(K).
\tag{M1}
\]

The coordinate action is the natural vertex permutation action. This
construction theorem is a substantial retained input; neither its
polytopality nor the identification of the full group in (M1) is
inferred solely from an order calculation.

\textbf{Corollary C.Parity.} Halving \(2^K\) at its vertex endpoint
gives the exact group

\[
E_m\rtimes\operatorname{Aut}(K),\qquad
E_m=\{x\in\mathbf F_2^m:\sum x_i=0\},
\tag{M2}
\]

of order \(2^{m-1}|\operatorname{Aut}(K)|\), acting as the full
type-preserving group of the regular hypertope in Corollary C.Halving.

\textbf{Proof.} The vertex graph of \(2^K\) is the m-cube, bipartitioned
by coordinate parity. Its color-preserving subgroup of (M1) is exactly
(M2), so Corollary C.Halving identifies the whole halving group.
Alternatively, let a flip the coordinate of a base vertex v of K and let
b be its first distinguished reflection, carrying v to an adjacent
vertex \(v'\). The product \(a'b=(aba)b\) is translation by
\(e_v+e_{v'}\). The unchanged reflections generate Aut(K). The original regular group is transitive on edges. The common graph-generation corollary in [U] therefore gives the whole augmentation module \(E_m\), and proves generation of the displayed semidirect product. \(\square\)

The source uses the dual convention \(\widehat P=(2^{K^*})^*\). Consequently m
in its formulas is the number of facets of K, and reversing all types
converts (M2) exactly to its construction.

\textbf{Corollary C.Polygon (polygon inputs).} For a p-gon K,
\(p\ge 3\), the result is a self-dual orientable regular map of type
\(\{p,p\}\), with

\[
V=F=2^{p-1},\qquad E=p2^{p-2},\qquad
|\operatorname{Aut}|=p2^p,\qquad
g=1+2^{p-3}(p-4).
\tag{M3}
\]

\textbf{Proof.} The two vertex colors give V and F after ordering the
map types as (color 0, old square, color 1). Each edge of the p-gon
specifies \(2^{p-2}\) coordinate squares, giving E. The fork diagram
becomes \{p,p\}. A single coordinate flip conjugates the two new
endpoint reflections and exchanges the colors, supplying the duality.
Extend the reflection-sign character of the dihedral Aut(K) trivially
over \(E_m\). Each of the three new distinguished reflections has sign
\(-1\), so the chamber orientation is well defined and reverses across
every adjacency. This proves orientability. Euler's formula
\(V-E+F=2-2g\) gives (M3). \(\square\)

For \(p=3\) the map is the tetrahedron: the group is
\(E_{3}\rtimes S_{3}\) of order 24 with the spherical \(A_{3}\) marking
and its universal Coxeter geometry. For \(p=4\) it is the square toroid
conventionally denoted \(\{4,4\}_{2,2}\). The following toroidal section proves that lattice
identification and identifies the original dual generalized cube with
\(\{4,4\}_{4,0}\) by explicit face-poset and marking maps.

\textbf{Corollary C.Spherical (all locally spherical inputs and
exceptional orders).} In output rank \(n\ge 4\), the locally spherical
cases of the source's dual generalized-cube construction are obtained from K
equal to the \((n-1)\)-simplex or \((n-1)-cube\), and additionally the
dodecahedron for \(n=4\) and the 120-cell for \(n=5\). Their types are
respectively spherical \(D_n\), affine \(\widetilde B_{n-1}\), and the
rank-four/rank-five compact-hyperbolic forks. The last two finite groups
have orders

\[
2^{11}\cdot120=245760,
\qquad 2^{119}\cdot14400.
\tag{M4}
\]

\textbf{Proof.} Put \(L=K^*\) to use the coordinate model. Each of its
two new vertex residues is L itself. If the output is locally spherical,
L must therefore be a universal spherical polytope. Retain the classical
finite string-Coxeter list. The lattice hypothesis excludes digonal
rank-two sections, so every relevant adjacent label is at least three.
In rank at least four a proper three-node fork with equal labels q,q is
spherical only for \(q=3\), where q is the first label of L. In rank
five, the remaining candidate \(F_{4}=\{3,4,3\}\) is excluded because
its proper four-node fork is affine. The finite list leaves exactly the
simplex, cross-polytope, icosahedron and 600-cell, the last two only in
their respective ranks. For rank at least six only the first two remain.

Conversely, all proper faces of these L are simplices. At an old j-face
of \(2^L\), \(j\ge 2\), let F be its associated \((j-1)\)-face of L. The
lower and upper sections vary independently. After halving, the maximal
residue is the product of the compressed cube \(2^F\) with the spherical
section L/F. As F has j vertices, the first factor has group
\(E_j\rtimes S_j\) and its \(D_j\) marking; its order \(2^{j-1}j!\)
equals that of the finite Coxeter group \(D_j\), proving equality with
the universal spherical geometry. Here \(D_{2}=A_{1}\times A_{1}\) and
\(D_{3}=A_{3}\). Vertex residues are L, so all maximal residues, and
consequently all proper residues, are genuinely spherical geometries. A
diagram alone is not used as a substitute for those residues.

The new global diagrams follow from Corollary C.Halving. The retained
finite, affine and compact-hyperbolic Coxeter classifications identify
them as stated. Formula (M2), with \(m=12\) or 120 and the classical
orders \(|\operatorname{Aut}(dodecahedron)|=120\) and
\(|\operatorname{Aut}(120-cell)|=14400\), gives (M4). For the simplex
\(m=n\) and \(|\operatorname{Aut}(K)|=n!\), giving \(D_n\) order
\(2^{n-1}n!\). For the \((n-1)-cube\), \(m=2(n-1)\) and
\(|\operatorname{Aut}(K)|=2^{n-1}(n-1)!\), giving \(2^{3n-4}(n-1)!\).
The finite hyperbolic examples are proper hypertopes because the fork
diagrams have a vertex of degree three. \(\square\)

The exact all-rank toroidal labels require more than these orders and
residues. They are established in the next section by an explicit
coordinate identification.

\section{Exact toroidal face posets and translation
lattices}\label{exact-toroidal-face-posets-and-translation-lattices}

Let \(r\ge 2\) and let L be the r-dimensional cross-polytope with
vertices labelled \(+e_i,-e_i\), \(1\le i\le r\). A proper face of L
chooses at most one vertex from each antipodal pair. Write \(Q_r\) for
the cubical tessellation of \(\mathbb R^r\) whose vertices are
\(\mathbb Z^r\) and whose cells are products of unit intervals and
integer singletons. The quotient \(Q_r/4\mathbb Z^r\) has a formal least
and greatest face adjoined when viewed as an abstract polytope. These
formal faces are also understood below.

For each pair of binary coordinates belonging to \(+e_i,-e_i\), use the
Gray map

\[
(0,0)\mapsto0,\quad(1,0)\mapsto1,\quad(1,1)\mapsto2,\quad(0,1)\mapsto3
\quad\text{in }\mathbb Z/4\mathbb Z.
\tag{T1}
\]

\textbf{Lemma C.Gray.} The product of these maps is a rank-preserving
face-poset isomorphism \(2^L\cong Q_r/4\mathbb Z^r\). Binary weight
modulo two becomes \(\sum_i x_i\) modulo two in the cubical
coordinates.

\textbf{Proof.} In Danzer's construction a proper face of \(2^L\) allows
the bits indexed by a face of L to vary, and fixes every other bit. Thus
in any antipodal pair either both bits are fixed or exactly one bit
varies. Under (T1), varying a single bit describes precisely one edge of
the four-cycle: its four possible edges arise from the two choices of
varying bit and the two values of the fixed bit. A face therefore
becomes a product of four-cycle vertices and edges. Conversely every
such product uniquely specifies the varying bit in each free pair and
the remaining fixed bits. Its dimension is the number of free pairs.
Inclusion is preserved and reflected coordinatewise, and formal faces
match. These products are exactly the cubical cells modulo
\(4\mathbb Z^r\), with their incidence. The four cases of (T1) have
binary parities 0,1,0,1 respectively, proving the last assertion.
\(\square\)

This is a face-poset identification, rather than an inference from
orders. It also handles the dual in the target source explicitly. For a
cubical face F with free coordinate set S, write its free intervals as
\([m_i,m_i+1]\) for \(i\in S\) and its fixed coordinates as \(m_j\) for
\(j\notin S\). The dual face \(F^*\) has fixed coordinate \(m_i+1/2\) for
\(i\in S\) and free interval \([m_j-1/2,m_j+1/2]\) for \(j\notin S\).
This reverses containment: whenever an endpoint is fixed in a subface of
F, the corresponding free interval of the dual contains the fixed
midpoint of \(F^*\). It is a bijection onto the shifted cubical
tessellation, with dimensions adding to r. Subtracting
\((1/2,\ldots ,1/2)\) identifies that tessellation with \(Q_r\). The
construction commutes with translations by \(4\mathbb Z^r\), and
consequently induces an order-reversing bijection on the quotient. Thus
for the r-cube \(K=L^*\), the source's polytope \((2^{K^*})^*\) is
indeed the cubical toroid \(\{4,3^{r-2},4\}_{(4,0,\ldots,0)}\). The
string is \(\{4,4\}\) when \(r=2\). Reversing the source marking moves its
last-end halving to first-end halving of \(2^L\). Type reversal is kept
explicit; no unproved compatibility of a chosen self-duality with
halving is needed.

\subsection{Markings and the affine
quotient}\label{markings-and-the-affine-quotient}

Use the following reflections of \(\mathbb R^r\):

\[
\begin{aligned}
a(x_1,x_2,\ldots,x_r)&=(1-x_1,x_2,\ldots,x_r),\\
b_i&:\ x_i\leftrightarrow x_{i+1}\quad(1\le i<r),\\
c(x_1,\ldots,x_{r-1},x_r)&=(x_1,\ldots,x_{r-1},-x_r).
\end{aligned}
\tag{T2}
\]

They generate \(W=\mathbb Z^r\rtimes B_r\), where \(B_r\) is the signed
permutation group. Indeed the \(b_i,c\) generate \(B_r\), including the
reflection \(x_1\mapsto -x_1\); its product with a is translation by
\(±e_1\). Signed permutations then supply every integer translation. The
simplex \(1/2\ge x_1\ge x_2\ge \cdots \ge x_r\ge 0\) is a fundamental
chamber. One way to verify this without a presentation theorem is to
translate each coordinate to \([-1/2,1/2]\), make it nonnegative, and
order the coordinates; a point with strict inequalities has a unique
such representative. Its reflected walls are the displayed generators.
Its vertices are the barycentres of a nested flag of cubical faces, so
its chamber system is the barycentric subdivision of \(Q_r\).

The marking in (T2) also agrees with Danzer's marking under (T1).
Flipping the \(+e_1\) bit induces \(x_1\mapsto 1-x_1\) modulo four;
interchanging antipodal pairs induces the \(b_i\); swapping the two bits
in the last pair induces \(x_r\mapsto -x_r\). These are the usual base
vertex/base flag generators of \(2^L\). Hence the isomorphism of Lemma
C.Gray identifies the marked action, not only the unmarked face poset.
The action of W on \(Q_r/4\mathbb Z^r\) has kernel exactly
\(\Lambda =4\mathbb Z^r\): an affine signed permutation acting trivially
on all vertices has translation zero modulo four, and its linear part is
the identity modulo four; testing the unit vectors forces that signed
permutation to be the identity.

Put \(b=b_1\) and replace a by

\[
d=aba:\ (x_1,x_2,x_3,\ldots,x_r)\mapsto(1-x_2,1-x_1,x_3,\ldots,x_r).
\tag{T3}
\]

Let \(E=\{v\in \mathbb Z^r\) : \(\sum v_i\) is even\}.

\textbf{Lemma C.Affine.} \(H=\langle d,b_1,\ldots ,b_{r-1},c\rangle\) is
exactly \(E\rtimes B_r\). It has index two in W and contains
\(\Lambda\). Its fundamental chamber is

\[
x_1\ge x_2\ge\cdots\ge x_r\ge0,\qquad x_1+x_2\le1.
\tag{T4}
\]

Its chamber system is the halved cubical geometry, and the marked finite
halving group and geometry are \(H/\Lambda\) and the quotient of that
chamber system by \(\Lambda\), respectively.

\textbf{Proof.} Every displayed generator preserves the parity of an
integral vertex's coordinate sum, so \(H\le E\rtimes B_r\). It contains
\(B_r\). The product of d with the reflection in \(B_r\)
\((x_1,x_2)\mapsto (-x_2,-x_1)\) is translation by \(e_1+e_2\).
Conjugating by \(B_r\) gives every \(±e_i±e_j\). These span E: subtract
multiples of \(e_i-e_r\), \(i<r\), to leave an even multiple of
\(e_r\), and use \(2e_r=(e_r+e_1)+(e_r-e_1)\). This works also for
\(r=2\). Thus equality holds and the index is two. Since every \(4e_i\)
is in E, \(\Lambda \le H\).

The region (T4) is the union of the preceding W chamber and its
reflection in the a-wall \(x_1=1/2\). On their common wall the two
chambers glue; the other wall \(x_1=x_2\) is paired with its a-image
\(x_1+x_2=1\). All other walls agree. Because \(W=H\sqcup Ha\) and the
interiors of W chambers are disjoint, these doubled chambers tile
\(\mathbb R^r\) under H, each once. Their wall reflections are exactly
\(d,b_1,\ldots ,b_{r-1},c\). Equivalently, their chamber labels are the
old cubical flags with the initial vertex chosen in one parity class,
pairing the two old flags differing only at that vertex. The two ends of
the old edge become the two new vertex types; the higher faces retain
their type. This is precisely the graph-halving construction in [U].

Finally, translations in \(\Lambda\) preserve both vertex colors and
take cubical faces to cubical faces. The same flag pairing therefore
commutes with \(\Lambda\). Every finite chamber has a lift, and two
lifts yield the same finite chamber exactly when they differ by
\(\Lambda\). Wall adjacencies are the marked reflections modulo
\(\Lambda\). The maximal parabolics are the images of the corresponding
wall subgroups, so the typed coset geometry also agrees. Regularity,
thinness and residual connectedness of the finite geometry follow from
the already proved halving theorem; they are not inferred solely from
being a group quotient. \(\square\)

For \(r\ge 3\), the walls (T4) have the affine fork diagram
\(\widetilde B_r\): d and \(b_1\) commute; both meet \(b_2\) with
product order three, successive \(b_i\) have order three, and
\(b_{r-1}c\) has order four. All other distinct wall pairs commute. This
follows directly from (T2)--(T3), or from their mirror angles. Thus its
rank-(r+1) quotient has the source's type \(\{4,3^{r-3},3/3\}\). Here
the quotient lattice is exactly \(4\mathbb Z^r\) in the original
coordinates; the subscript \((4,0,\ldots ,0)\) denotes the lattice
generated by that vector's signed-permutation orbit. This proves the
source's all-rank toroidal identification, including the marking. Its
group order is \([E:4\mathbb Z^r]|B_r|=(4^r/2)2^r\) \(r!=2^{3r-1}r!\),
agreeing with the separate parity-module calculation. The order is a
check, not the proof.

The coordinate convention can be checked directly against Ens's primary
paper: its equations (3.1), (4.1) are (T2), (T3) for \(r=3\); its §2
definition of \(\Lambda _s\) uses the signed-permutation orbit of s, and
Theorem 4.3 denotes the corresponding Coxeter-complex quotient by that
lattice. Thus \(s=(4,0,0)\) here is literally the same parameter. The
target extends this convention to larger r. No classification theorem
about all invariant lattices is needed. See {[}Ens, pp.70--75{]}. Only
that lattice convention is borrowed here, not its classification
theorem.

\subsection{The rank-three square map}\label{the-rank-three-square-map}

When \(r=2\), reorder the halved types as (even vertex, old square, odd
vertex). An old unit square joins its two even corners as an edge of the
new map; the four even neighbors of an old odd vertex bound a new square
face. This describes the ordinary square tessellation with vertex
lattice E and edge directions

\[
f_1=(1,1),\qquad f_2=(-1,1).
\tag{T5}
\]

Every even vertex has four such incident edges, every displayed
quadrilateral is a face, and the incidences are exactly those of
halving. Thus this is an isomorphism of maps, not a comparison of their
counts. In this basis

\[
4e_1=2f_1-2f_2,\qquad4e_2=2f_1+2f_2.
\]

The quotient lattice has basis \((2,-2),(2,2)\) in square coordinates.
Negating the first vector gives the customary pair \((-2,2),(2,2)\), so
the map is \(\{4,4\}_{2,2}\). Lemma C.Gray and the explicit duality
already identify the original map as \(\{4,4\}_{4,0}\). This proves both
\(p=4\) identifications in [MW20].

\section{Scope, rank qualification and attribution}
The specified edition is arXiv:2001.11082v1, dated 29 January 2020,
eleven physical pages. Its generalized-cube construction, polygon
applications, complete locally spherical list and exact toroidal
identifications are the conclusions recovered above. The printed
names of toroidal geometries are justified by face-poset and marked
lattice maps, not just group orders.

There is a rank qualification in Remark 4.1. Retain the Pellicer
extension construction used in that remark. After endpoint-$2s$
halving its three-node fork triangle has labels $(p,p,s)$. In output
rank at least four, with finite $p,s\ge3$, this is a proper rank-three
diagram and is not finite, since $2/p+1/s\le1$. The result is therefore
not locally spherical. In rank three all proper residues have rank at
most two, and these finite parameters give finite dihedral residues.
The unqualified wording consequently does not apply in rank three
under the source's definition. No classification of the entire
extension construction is asserted here.

The comparison against every example in the source's reference [9],
and its unused alternating-semiregular identification, are excluded
auxiliary assertions. Danzer's construction and full original
automorphism group, classical Coxeter classifications and convex
polytope orders, and standard orientation and Euler-characteristic
facts are retained inputs. The selected pages of Ens fix the toroidal
lattice convention; its classification theorem is not used to prove
our particular quotient identifications.

The new use of the common binary module theorem is exact: conjugates
of one edge difference generate its augmentation module, so the same
generation argument serves this application and the equal-parameter
polyhedra. It does not supply an independent proof of Danzer's
polytopality theorem. The present application remains written rather
than fully Lean-formalized; the common endpoint file verifies only its
specified fixed-overlap identity.

\section*{References}
{[}U{]} \emph{Changes of chamber generators and binary translation
groups}, the common paper of Topic C.

{[}MW20{]} A. Montero and A. Ivi\'c Weiss,
\emph{Locally spherical hypertopes from generalised cubes},
arXiv:2001.11082v1, 29 January 2020.

{[}Danzer{]} L. Danzer, \emph{Regular incidence-complexes and
dimensionally unbounded sequences of such. I}, North-Holland
Mathematics Studies 87 (1984), 115--127.

{[}Ens{]} E. Ens, \emph{Rank 4 toroidal hypertopes}, Ars Mathematica
Contemporanea 15 (2018), 67--79; selected pages 70--75, especially
equations (3.1), (4.1) and the convention before Theorem 4.3.

{[}MS02{]} P. McMullen and E. Schulte, \emph{Abstract Regular
Polytopes}, Cambridge University Press, 2002.
\end{document}
